\section{布尔代数与逻辑化简}
\label{sec:04-Discrete-Mathematics/07-Boolean-Algebra-and-Automata/01}

你接手了一段访问控制规则，前一个同事是这么写的：

\begin{quote}
用户已登录，并且是管理员或资源所有者；或者请求来自紧急恢复通道，但账号没有被冻结。
\end{quote}

把它翻译成变量：

\[
A=\text{已登录},\quad
B=\text{是管理员},\quad
C=\text{是资源所有者},\quad
D=\text{来自紧急恢复通道},\quad
E=\text{账号被冻结}.
\]

写成公式：

\[
F = \big(A\land(B\lor C)\big)\ \lor\ (D\land\neg E).
\]

现在另一个开发者重构了这段逻辑，写成了不同的表达式——格式看起来不一样，代码审查者拿不准两版到底等不等价。直觉上好像一样，但之前就因为逻辑等价性搞错过一次线上事故。于是你打开真值表。

五个变量，\(2^5=32\) 行。逐行对一遍，两个公式每行都一样，放心了。接下来系统升级，变量增加到八个，\(2^8=256\) 行，还能应付。但产品说下个版本要加三个条件——11 个变量，2048 行。更麻烦的是，你想把电路实现的门数从 8 个降到 4 个，靠肉眼在真值表里找化简可太难了。

你手上的不是数学题，是一个会用真值表、也快被真值表逼疯的工程师处境。这就是布尔代数的入口：我们需要一套代数运算，能像解方程一样对逻辑表达式做等价变形，而不是每次重新比对所有行。

\subsection{两个值与三种基本运算的算术面孔}

用 \(1\) 表示真、\(0\) 表示假。对 \(x,y\in\set{0,1}\)，定义

\begin{itemize}
\tightlist
\item
  合取 \(x\land y\)：二者都为 \(1\) 时为 \(1\)；
\item
  析取 \(x\lor y\)：至少一个为 \(1\) 时为 \(1\)；
\item
  否定 \(\neg x\)：交换 \(0\) 与 \(1\)。
\end{itemize}

这三个运算恰好能用普通算术写出来：

\[
x\land y=xy,
\qquad
\neg x=1-x,
\qquad
x\lor y=x+y-xy.
\]

第一眼会觉得这很妙——不需要新符号，直接用加减乘就能算逻辑。但算一个具体值立刻暴露陷阱：\(1\lor1=1\)，而按上面公式 \(1+1-1\cdot1=1\)，确实对。可如果把 \(\lor\) 直接当成普通加法，\(1+1=2\) 就跑出 \(\{0,1\}\) 了。布尔代数里的``加''指 OR，``乘''指 AND——名字借用算术，规矩却是自己的。尤其是幂等性 \(x\lor x=x\) 和 \(x\land x=x\)，在普通代数里横竖不对。

不过把算术公式当检查工具确实好用。比如验证 De Morgan 时，左边 \(\neg(x\land y)=1-xy\)，右边 \((\neg x)\lor(\neg y)=(1-x)+(1-y)-(1-x)(1-y)=1-x+1-y-(1-x-y+xy)=1-xy\)，两边相等。算术没有代替逻辑，它只是给逻辑搭了座桥。

集合提供了另一套平行直觉。固定全集 \(U\)，把命题 \(x\) 对应为使它为真的对象集合 \(X\subseteq U\)：

\[
x\land y\longleftrightarrow X\cap Y,
\qquad
x\lor y\longleftrightarrow X\cup Y,
\qquad
\neg x\longleftrightarrow U\setminus X.
\]

这就是布尔代数的甜头：逻辑公式、开关电路、集合运算共享同一套代数结构。你在一个领域练熟的变形，搬运到另一个领域照样合法。

\subsection{运算律：哪些变形被允许}

布尔代数的运算律和普通代数有重叠，但分叉的地方恰好是关键。

首先是边界元：

\[
x\lor0=x,\qquad x\land1=x,
\]
\[
x\lor1=1,\qquad x\land0=0.
\]

0 对 OR 像是加法零元，但对 AND 却是湮灭元——这跟普通乘法不同，普通乘法里 0 才是零元，1 是单位元。

布尔代数最不合群的两条是幂等和互补：

\[
x\lor x=x,\qquad x\land x=x,
\]
\[
x\lor\neg x=1,\qquad x\land\neg x=0.
\]

幂等性意味着重复同一个条件不产生额外信息，互补性则是二值逻辑的根基：一件事要么真要么假，没有中间态。

AND 与 OR 各自满足交换律和结合律，分配律更是双向的：

\[
x\land(y\lor z)=(x\land y)\lor(x\land z),
\]
\[
x\lor(y\land z)=(x\lor y)\land(x\lor z).
\]

第二条在普通代数里不存在——\(x+(y\times z)\) 不等于 \((x+y)\times(x+z)\)。布尔代数里 AND 和 OR 的地位完全对称，这就是对偶原理的根源：任意一条恒等式同时交换 \(\land\leftrightarrow\lor\) 与 \(0\leftrightarrow1\)，得到的仍是一条有效恒等式。记住一半，对偶自动给出另一半。

De Morgan 律把否定推进括号，在工程文档里被重复发现无数次：

\[
\neg(x\land y)=(\neg x)\lor(\neg y),
\]
\[
\neg(x\lor y)=(\neg x)\land(\neg y).
\]

读回自然语言：``并非两人都满足条件''等于``至少一人不满足''；``并非两人中有一个满足''等于``两人都不满足''。多数逻辑漏洞恰好出在否定范围上面——人脑对``不是 A 并且 B''这种嵌套否定天然不擅长，De Morgan 律是代数化的防错机制。

\subsection{吸收律：多余的零件自行消失}

考虑这一段表达式：

\[
x\lor(x\land y).
\]

逐个情况想：若 \(x=1\)，整个式子已为真，\(y\) 纯属多余；若 \(x=0\)，第二项 \(x\land y=0\) 也不可能救回来。因此

\[
x\lor(x\land y)=x.
\]

对偶地：

\[
x\land(x\lor y)=x.
\]

这就是吸收律。它不是机械展开然后碰巧消掉——它有直接的逻辑含义：一旦宽松条件 \(x\) 已覆盖所有成立情况，再搭一个更严的条件 \((x\land y)\) 并不会筛掉更多 case，当然也不会新增通过者。电路上，这意味着一个 AND 门和一个 OR 门整体可以删掉，因为它们的结果永远等于其中一根输入线。

再看一个经典的化简：

\[
(x\land y)\lor(x\land\neg y).
\]

把 \(x\) 提出来：

\[
x\land(y\lor\neg y)=x\land1=x.
\]

两项分别覆盖了 \(y\) 为真和为假的全部可能，\(y\) 其实没有施加任何约束。电路上，两个 AND 门、一个 OR 门和一个 NOT 门被缩成一根导线。这不止省了几个晶体管——它揭示了 \(y\) 在规则里是冗余变量：看上去它参与了决策，实际上它怎么取值都改变不了输出。一条规则也许能删掉整个判断条件，这在合规审计和代码简化中都有实际价值。

\subsection{所有布尔函数都能写成范式}

任意 \(n\) 元布尔函数都能用 AND、OR、NOT 表示，构造方式直接来自真值表：对每一行输出为 \(1\) 的赋值，写一个只在该行成立的合取项，然后把这些项全部析取起来。

比如函数 \(f(x,y,z)\) 只在

\[
(x,y,z)=(1,0,1)\quad\text{或}\quad(0,1,1)
\]

时为真。对应的两个合取项自然就是

\[
f=(x\land\neg y\land z)\lor(\neg x\land y\land z).
\]

这种``若干合取项再析取''的形式就是析取范式（DNF）。再提一下公因子：

\[
f=z\land\big((x\land\neg y)\lor(\neg x\land y)\big)=z\land(x\oplus y),
\]

其中异或 \(x\oplus y\) 表示恰有一个为真。从 DNF 出发提公因式化简，就是我们在访问控制场景里想做的事。

对称地，对每一行输出为 \(0\) 的赋值写一个只在该行失败的析取子句，再把这些子句全部合取，得到合取范式（CNF）。DNF 回答``哪些情况足以使它为真''，CNF 回答``必须同时满足哪些约束''。同一个函数，两种视角。

这里藏了一个务实边界：规范 DNF 最坏可含 \(2^n\) 个项，每项含 \(n\) 个文字。把公式转成标准形本身就可能指数膨胀。``任何逻辑都能写''跟``能写得短、算得快''是两件不同的事——存在性定理不承诺可行性。

\subsection{NAND 一个门就能干所有活}

既然 AND、OR、NOT 能表示任意布尔函数，它们构成功能完备的连接词集合。但更极端的结论是：单独一个 NAND

\[
x\uparrow y=\neg(x\land y)
\]

就足够了，因为

\[
\neg x=x\uparrow x,
\]
\[
x\land y=(x\uparrow y)\uparrow(x\uparrow y),
\]

再配合 De Morgan 律就能生成 OR。NOR 也同样完备。

这意味着硬件只需要生产一种标准门就能搭出所有组合逻辑。不过``只用一种门''不等于电路更省：只用 NAND 搭一个异或门，可能比直接用 AND/OR/NOT 多出一堆门。实际芯片设计里会在门数量、逻辑深度、扇入、功耗和延迟之间做多目标权衡。代数上文字更少的公式，不必然对应具体工艺下延迟最低的网表——逻辑等价不等于物理等价。

异或在奇偶校验与加法器中出现得特别频繁：

\[
x\oplus y=(x\land\neg y)\lor(\neg x\land y),
\]

它满足交换律与结合律，并且 \(x\oplus x=0\)。多个比特异或为 \(1\) 当且仅当其中有奇数个 \(1\)，这正是\hyperref[sec:04-Discrete-Mathematics/05-Coding-and-Polynomial-Methods/02]{``码距决定检测与纠错能力''}中奇偶校验的代数——同一个异或运算，一边在说算术加法不进位，一边在说码字差错检测。

\subsection{化简有多难：局部可做，全局很难}

小公式用恒等式手工化简。变量五六个的时候，Karnaugh 图把相邻真值格合并成矩形块，肉眼能直接读出某个变量在哪些区域可以消去。当年设计人员真的在方格纸上画卡诺图，那是一种把几何直觉注入逻辑设计的办法。

变量再多就得靠算法：Quine--McCluskey、二元决策图（BDD）、或直接把等价问题扔给 SAT 求解器。

判断两个 \(n\) 元公式是否等价，可以检查

\[
f\oplus g
\]

是否恒为 \(0\)——把等价问题转化成了可满足性问题。一般布尔公式的可满足性是 NP 完全的：不存在已知方法能对所有公式快速给出全局最优化简。局部恒等式每次化简都有效，但不保证最终到达门数最少的表达。

BDD 的例子尤其有教育意义：同一函数，在好的变量顺序下 BDD 很小，坏的顺序下可能指数大。表示方式的选择本身就是算法难度的一部分——不能只看最终公式有几行，还得看找到它的过程花了多少搜索。

实践中一般不追求绝对的全局最优。工程方法是：局部化简加启发式搜索，得到一个``足够好''的公式，剩下的留给综合工具去管。数学给出最优性的困难证明，恰是工程接受``足够好''的底气。

\subsection{二值逻辑装不下的真实系统}

布尔代数假设每个变量非真即假，并且运算是瞬时、确定的。走出芯片内部，这套假设就开始漏水了。

数据库 SQL 有 NULL，形成三值逻辑：\(x\lor\neg x\) 在 \(x\) 未知时不返回真，因为``未知 OR 未知的否定''不等于确定真。传感器判断带置信度，需要概率；数字电路有传播延迟，两个逻辑等价的表达式可能产生不同的瞬态毛刺——\(A\land B\) 和 \(B\land A\) 在数学上相等，但在物理走线上，谁先到、谁后到会决定输出是否出现一个纳秒级的错误脉冲。

访问控制还有优先级、默认拒绝和副作用。经典恒等式说你可以自由重排，但如果每个条件调用了一个带副作用的函数，重排可能改变执行结果。化简前必须确认：你化简的是纯二值逻辑，还是一个披着逻辑外衣的编程语言表达式。

布尔代数刻画的是纯二值函数的等价。它不涵盖时序电路的行为等价，不涵盖连续值推断，不涵盖带副作用或异常控制的程序语义。但这些边界恰好说明了它的精确适用范围——在它的疆域内，每一条等式都无条件成立。知道哪一步走出去就不再安全，才是真正懂了这套代数。
